\documentclass{article}

% Language and encoding
\usepackage[portuguese]{babel}
\usepackage[utf8]{inputenc}
\usepackage[T1]{fontenc}

% Page layout
\usepackage[a4paper,left=25mm,right=25mm,top=25mm,bottom=25mm,headheight=6mm,footskip=12mm]{geometry}
\setlength{\parindent}{1em}
\setlength{\parskip}{1ex}

% Core packages
\usepackage{amsmath}
\usepackage{booktabs}
\usepackage{indentfirst}
\usepackage{ragged2e}
\usepackage{setspace}
\usepackage{caption}
\usepackage{float}
\usepackage{longtable}
\usepackage{multirow}
\usepackage{graphicx}
\usepackage{adjustbox}
\usepackage{lmodern}
\usepackage{xcolor}
\usepackage{listings}
\usepackage{lastpage}
\usepackage{fancyhdr}
\usepackage{hyperref}

%% headers & footers
\fancyhf{}
\rhead{\small{\emph{ATS - Relatorio Geral}}}
\rfoot{\small{\thepage\ / \pageref{LastPage}}}
\pagestyle{fancy}
\renewcommand{\headrulewidth}{0.4pt}
\renewcommand{\footrulewidth}{0.4pt}

%% colors
\definecolor{engineering}{rgb}{0.549,0.176,0.098}
\definecolor{cloudwhite}{cmyk}{0,0,0,0.025}

%% source-code listings
\lstset{
 language=bash,
 basicstyle=\footnotesize\ttfamily,
 keywordstyle=\bfseries,
 numbers=left,
 numberstyle=\scriptsize\texttt,
 stepnumber=1,
 numbersep=8pt,
 frame=tb,
 float=htb,
 aboveskip=8mm,
 belowskip=4mm,
 backgroundcolor=\color{cloudwhite},
 showspaces=false,
 showstringspaces=false,
 showtabs=false,
 tabsize=2,
 captionpos=b,
 belowcaptionskip=12pt,
 breaklines=true,
 breakatwhitespace=false,
 commentstyle=\color{gray},
 lineskip=3pt,
 extendedchars=true
}

%% hyperreferences
\hypersetup{
 plainpages=false,
 pdfpagelayout=SinglePage,
 bookmarksopen=false,
 bookmarksnumbered=true,
 breaklinks=true,
 linktocpage,
 colorlinks=true,
 linkcolor=engineering,
 urlcolor=engineering,
 filecolor=engineering,
 citecolor=engineering,
 allcolors=engineering
}

%% path to figures
\graphicspath{{figures/}}

%% macros
\newcommand{\school}{Universidade do Minho}
\newcommand{\degree}{Licenciatura em Engenharia Informatica}
\newcommand{\projtitle}{Analise e Teste de Software}
\newcommand{\subtitle}{2025/2026}
\newcommand{\projauthor}{Trabalho Pratico}
\newcommand{\projname}{Relatorio Geral - Projetos 1 e 2}

\newcommand{\atsfigure}[3]{
\begin{figure}[H]
\centering
\IfFileExists{#1}{
    \includegraphics[width=#2\textwidth]{#1}
}{
    \fbox{\parbox[c][50mm][c]{0.9\textwidth}{\centering Placeholder para print: #1}}
}
\caption{#3}
\end{figure}
}

\renewcommand{\thesection}{\arabic{section}}
\setcounter{secnumdepth}{3}

\begin{document}

% Capa
\begin{titlepage}
\centering

\vspace{-15mm}
{\large \textbf{\textsc{\school}}}\\

\vspace{0.5cm}
{\large \degree}\\[8mm]

\vspace{2cm}
\IfFileExists{UM.png}{\includegraphics[width=52mm]{UM.png}}{\vspace{20mm}}

\vspace{1cm}
{\Large \projtitle}\\[2mm]
{\large \subtitle}\\[20mm]

{\Large \textbf{\projauthor}}\\
\vspace{0.5cm}
{\Large \projname}

\vspace{4cm}

{\large Francisco Barros A106894}\\[2mm]
{\large Tiago Martins A106927}\\[2mm]

\vspace{2cm}
{\Large \textbf{Abril 2026}}

\end{titlepage}

\renewcommand{\contentsname}{Indice}
\tableofcontents
\newpage
\listoffigures
\newpage
\listoftables
\newpage

\section{Introducao}
Este documento apresenta uma versao geral e consolidada do trabalho de Analise e Teste de Software aplicado ao \textbf{Projeto 1 (SpotifyUM)} e ao \textbf{Projeto 2 (SpotifUM)}.

A consolidacao foi feita com tres objetivos centrais:
\begin{itemize}
    \item preservar a rastreabilidade tecnica de todas as decisoes de teste;
    \item comparar de forma objetiva os resultados obtidos em duas toolchains diferentes;
    \item fornecer um artefacto unico, completo e reprodutivel para avaliacao.
\end{itemize}

A estrategia global combinou:
\begin{itemize}
    \item reforco manual de testes unitarios;
    \item geracao automatica de testes (EvoSuite e apoio por IA);
    \item analise estrutural por cobertura;
    \item analise de robustez da suite por mutacao;
    \item property-based testing;
    \item automacao de pipeline e repetibilidade de execucao.
\end{itemize}

\section{Contexto dos Projetos}
\subsection{Projeto 1 - SpotifyUM (Maven)}
O Projeto 1 assenta em Maven e segue uma organizacao classica de codigo Java com separacao de camadas de \texttt{Model}, \texttt{Controller}, \texttt{View}, \texttt{Utils} e excecoes.

Do ponto de vista de teste, o foco foi:
\begin{itemize}
    \item consolidar a suite JUnit existente;
    \item aumentar cobertura em utilitarios e fluxos criticos de controlador;
    \item integrar EvoSuite por perfis Maven.
\end{itemize}

\subsection{Projeto 2 - SpotifUM (Gradle)}
O Projeto 2 assenta em Gradle e usa uma pipeline de verificacao mais abrangente, incluindo JaCoCo, PIT e jqwik.

Do ponto de vista de teste, o foco foi:
\begin{itemize}
    \item reforco de testes em utilitarios e regras de negocio;
    \item medicao sistematica de cobertura;
    \item uso de mutacao para localizar fragilidades da suite;
    \item automacao integral em tarefas Gradle e script unico.
\end{itemize}

\section{Objetivos ATS e Criterios de Qualidade}
\subsection{Objetivos obrigatorios}
\begin{itemize}
    \item complementar e estabilizar suites de testes;
    \item usar geracao automatica de testes;
    \item medir cobertura e qualidade;
    \item aplicar mutacao e interpretar sobreviventes;
    \item aplicar property-based testing;
    \item apresentar relatorio tecnico com analise critica.
\end{itemize}

\subsection{Objetivos de excelencia}
\begin{itemize}
    \item garantir reproducibilidade de execucao por comandos simples;
    \item justificar trade-offs de tempo e qualidade;
    \item registar limitacoes externas sem perder validade metodologica;
    \item integrar assistencia por IA com revisao manual obrigatoria.
\end{itemize}

\section{Porque manter Maven num projeto e Gradle no outro}
A coexistencia de Maven e Gradle foi uma decisao deliberada para explorar ecossistemas complementares:
\begin{itemize}
    \item \textbf{Maven} no Projeto 1, por favorecer configuracao declarativa e separacao por \textit{profiles}, util para cenarios com modos de execucao distintos (suite normal, runtime EvoSuite, geracao EvoSuite).
    \item \textbf{Gradle} no Projeto 2, por facilitar composicao de tarefas ATS e automacao de pipeline end-to-end, incluindo relatorios em cadeia.
\end{itemize}

Esta abordagem permitiu comparar praticas em contextos reais de manutencao: previsibilidade de ciclo de vida (Maven) versus flexibilidade de orquestracao (Gradle).

\section{Estrategia Transversal de Teste}
\subsection{JUnit}
JUnit foi a base comum dos dois projetos, com enfase em:
\begin{itemize}
    \item cenarios de fronteira;
    \item validacao de excecoes e caminhos negativos;
    \item garantia de regressao funcional apos alteracoes.
\end{itemize}

\subsection{EvoSuite}
EvoSuite foi integrado nos dois projetos para geracao automatica de testes. Para manter iteracoes curtas e reduzir custo computacional, a geracao foi limitada a 3 classes por projeto.

\subsection{JaCoCo e PIT}
No Projeto 2, JaCoCo e PIT foram usados em conjunto:
\begin{itemize}
    \item JaCoCo para visibilidade estrutural (o que e executado);
    \item PIT para eficacia semantica dos testes (o que realmente deteta defeitos).
\end{itemize}

\subsection{Property-Based Testing (jqwik)}
No Projeto 2, jqwik foi aplicado em propriedades de invariantes para aumentar diversidade de entradas e detetar casos nao previstos por testes exemplo-a-exemplo.

\section{Projeto 1 - SpotifyUM (Maven)}
\subsection{Intervencao principal}
O Projeto 1 recebeu reforco em testes utilitarios e de controlador, com enfase nos seguintes pontos:
\begin{itemize}
    \item serializacao (round-trip e erro de IO);
    \item validacoes de recursos multimedia e URL;
    \item autenticacao e fluxos de registo/login;
    \item persistencia/import-export em cenarios com ficheiros temporarios.
\end{itemize}

\subsection{Estado do build e testes observados}
\begin{table}[H]
\centering
\begin{tabular}{p{0.58\textwidth}c}
\toprule
Verificacao pratica & Resultado \\
\midrule
\texttt{mvn -DskipTests compile} & Sucesso (exit 0) \\
\texttt{mvn test} sem perfil adicional & Falha por dependencias de teste/EvoSuite runtime \\
\texttt{mvn -Ptests-junit5,evosuite-runtime test -DskipTests} & Sucesso de compilacao (exit 0) \\
\bottomrule
\end{tabular}
\caption{Verificacao de usabilidade no Projeto 1}
\end{table}

\subsection{Perfis Maven relevantes}
A configuracao Maven foi estruturada com perfis dedicados:
\begin{itemize}
    \item \texttt{tests-junit5}: dependencias JUnit 5 e Mockito;
    \item \texttt{evosuite-generate}: plugin de geracao de testes EvoSuite;
    \item \texttt{evosuite-runtime}: runtime necessario para classes \texttt{\_ESTest}.
\end{itemize}

\subsection{EvoSuite limitado a 3 classes}
No Projeto 1, o parametro \texttt{evosuite.cuts} foi restringido para:
\begin{itemize}
    \item \texttt{org.Model.SpotifUM}
    \item \texttt{org.Model.Album.Album}
    \item \texttt{org.Model.Music.Music}
\end{itemize}

\subsection{Justificacao tecnica desta limitacao}
A geracao automatica em todas as classes de dominio tem custo elevado e pode alongar excessivamente cada iteracao. A limitacao para 3 classes permite:
\begin{itemize}
    \item demonstrar integracao funcional da ferramenta;
    \item validar o processo geracao/exportacao/execucao;
    \item libertar tempo para analise de mutantes, cobertura e consolidacao de testes manuais.
\end{itemize}

\section{Projeto 2 - SpotifUM (Gradle)}
\subsection{Intervencao principal}
O Projeto 2 foi tratado como pipeline ATS completo, incluindo:
\begin{itemize}
    \item \texttt{test} com JUnit 5;
    \item \texttt{jacocoTestReport} para cobertura;
    \item \texttt{pitest} para mutacao;
    \item \texttt{evosuiteGenerate} para geracao automatica;
    \item \texttt{atsFullPipeline} para execucao agregada.
\end{itemize}

\subsection{Estado do build e testes observados}
\begin{table}[H]
\centering
\begin{tabular}{p{0.58\textwidth}c}
\toprule
Verificacao pratica & Resultado \\
\midrule
\texttt{bash ./gradlew testClasses} & Sucesso (exit 0) \\
\texttt{bash ./gradlew test} & Sucesso (exit 0) \\
\texttt{./gradlew test} sem permissao de execucao no wrapper & Falha de permissao no ambiente atual \\
\bottomrule
\end{tabular}
\caption{Verificacao de usabilidade no Projeto 2}
\end{table}

\subsection{EvoSuite limitado a 3 classes}
No Projeto 2, a lista \texttt{evosuiteTargets} foi restringida para:
\begin{itemize}
    \item \texttt{org.spotifumtp37.util.SongTypeAdapter}
    \item \texttt{org.spotifumtp37.util.SubscriptionPlanAdapter}
    \item \texttt{org.spotifumtp37.util.LocalDateTimeAdapter}
\end{itemize}

\subsection{Limitacao conhecida do EvoSuite neste projeto}
A tarefa de EvoSuite usa container Docker e pode encontrar limitacoes de compatibilidade de bytecode dependendo da imagem/runtime em uso, especialmente em projetos com Java mais recente. Esta e uma limitacao externa de toolchain e nao invalida a integracao tecnica nem o racional de uso.

\subsection{Matriz de resultados (iteracao reportada)}
\begin{table}[H]
\centering
\begin{tabular}{lr}
\toprule
Metrica & Valor \\
\midrule
Total de testes JUnit & 141 \\
Taxa de sucesso JUnit & 100\% \\
Line coverage (JaCoCo) & 40.02\% \\
Branch coverage (JaCoCo) & 43.69\% \\
Muta\c{c}oes geradas (PIT) & 890 \\
Mutation score (PIT) & 30\% \\
No coverage (PIT) & 587 \\
Test strength (PIT) & 87\% \\
\bottomrule
\end{tabular}
\caption{Resumo de metricas consolidadas do Projeto 2}
\end{table}

\atsfigure{figures/junit_report_overview.png}{0.92}{Resumo da execucao JUnit no Projeto 2}
\atsfigure{figures/jacoco_overview.png}{0.92}{Resumo de cobertura JaCoCo no Projeto 2}
\atsfigure{figures/pit_overview.png}{0.92}{Resumo de mutacao PIT no Projeto 2}

\section{Geracao automatica de testes com assistencia por IA}
\subsection{Porque usar IA neste contexto}
A assistencia por IA foi usada como amplificador de produtividade, sobretudo para:
\begin{itemize}
    \item propor listas de casos-limite e classes de equivalencia;
    \item acelerar criacao de esqueletos de testes unitarios;
    \item sugerir cenarios negativos e casos de excecao;
    \item sugerir propriedades para testes com jqwik.
\end{itemize}

\subsection{Controlo de qualidade aplicado}
Para garantir rigor academico e tecnico, todo output de IA foi sujeito a:
\begin{itemize}
    \item revisao manual linha-a-linha;
    \item ajuste ao comportamento real do codigo;
    \item execucao local obrigatoria antes de aceitar qualquer teste;
    \item rejeicao de casos redundantes, inconsistentes ou sem valor semantico.
\end{itemize}

\subsection{Posicionamento metodologico}
A IA nao substituiu engenharia de teste. Foi usada para reduzir custo mecanico e libertar tempo para atividades de maior valor:
\begin{itemize}
    \item analise de sobreviventes PIT;
    \item melhoria de cobertura em zonas de risco;
    \item interpretacao critica de resultados.
\end{itemize}

\section{Comparacao consolidada Projeto 1 vs Projeto 2}
\begin{longtable}{>{\raggedright\arraybackslash}p{0.26\textwidth} >{\raggedright\arraybackslash}p{0.33\textwidth} >{\raggedright\arraybackslash}p{0.33\textwidth}}
\toprule
\textbf{Dimensao} & \textbf{Projeto 1 (Maven)} & \textbf{Projeto 2 (Gradle)} \\
\midrule
Sistema de build & Ciclo declarativo + perfis & Tarefas composiveis + automacao agregada \\
Execucao base de testes & Sensivel a perfis ativos & Direta com \texttt{bash ./gradlew test} \\
EvoSuite & Integrado por perfil + runtime dedicado & Integrado por task Docker \\
Escopo EvoSuite & 3 classes alvo & 3 classes alvo \\
Cobertura estrutural & Nao centralizada no fluxo atual & JaCoCo integrado no fluxo principal \\
Mutacao & Nao foco principal nesta iteracao & PIT integrado e operacional \\
PBT & Nao foco principal nesta iteracao & jqwik aplicado em propriedades \\
Automacao end-to-end & Parcial & Completa via task/script \\
Risco operacional dominante & Gestao de perfis e runtime de teste & Compatibilidade de runtime EvoSuite + permissao wrapper \\
\bottomrule
\end{longtable}

\section{Reprodutibilidade}
\subsection{Comandos de referencia - Projeto 1}
\begin{lstlisting}[caption={Comandos chave de reproducao - Projeto 1},label={lst:p1-commands}]
cd Projeto1/SpotifyUM
mvn -DskipTests compile
mvn -Ptests-junit5,evosuite-runtime test
mvn -Pevosuite-generate test
\end{lstlisting}

\subsection{Comandos de referencia - Projeto 2}
\begin{lstlisting}[caption={Comandos chave de reproducao - Projeto 2},label={lst:p2-commands}]
cd Projeto2
bash ./gradlew test
bash ./gradlew jacocoTestReport pitest
bash ./gradlew evosuiteGenerate
bash ./gradlew atsFullPipeline
\end{lstlisting}

\subsection{Pipeline unica de verificacao}
\begin{lstlisting}[caption={Pipeline ATS agregada no Projeto 2},label={lst:pipeline-ats}]
cd Projeto2
bash ./gradlew atsFullPipeline --no-daemon
\end{lstlisting}

\section{Riscos, Limitacoes e Mitigacoes}
\subsection{Riscos identificados}
\begin{itemize}
    \item dependencia de perfis Maven para compilar/executar todas as variantes de testes no Projeto 1;
    \item dependencia de runtime/container para EvoSuite no Projeto 2;
    \item cobertura e mutation score ainda com margem de melhoria em zonas menos testadas.
\end{itemize}

\subsection{Mitigacoes adotadas}
\begin{itemize}
    \item comandos documentados por perfil e por etapa;
    \item limitacao de escopo EvoSuite para iteracoes mais curtas;
    \item reforco incremental orientado por metricas (JaCoCo/PIT);
    \item uso de IA com validacao humana para acelerar sem perder qualidade.
\end{itemize}

\section{Plano de melhoria para iteracoes futuras}
\begin{itemize}
    \item aumentar cobertura em classes de negocio com maior no-coverage no PIT;
    \item reduzir mutantes sobreviventes com testes dirigidos por operador de mutacao;
    \item alargar property-based testing para mais modulos e invariantes de dominio;
    \item estabilizar pipeline EvoSuite em ambiente controlado de versoes Java;
    \item manter verificacao de regressao apos cada refatoracao relevante.
\end{itemize}

\section{Conclusao}
O relatorio geral evidencia uma evolucao tecnica consistente nos dois projetos, com abordagens complementares:
\begin{itemize}
    \item no Projeto 1, consolidacao da suite e integracao por perfis Maven;
    \item no Projeto 2, pipeline ATS completa com cobertura, mutacao e automacao.
\end{itemize}

A decisao de aplicar EvoSuite a apenas 3 classes por projeto foi deliberada e adequada ao prazo, permitindo equilibrio entre profundidade tecnica e tempo de execucao.

A assistencia por IA contribuiu para acelerar geracao e refinamento de testes, mantendo rigor atraves de revisao manual e validacao por execucao local. Com isso, o processo resultou em maior produtividade sem comprometer qualidade metodologica.

Em sintese, os dois projetos encontram-se com base solida para avaliacao ATS, com evidencias reproduciveis, comparacao tecnica clara e roteiro concreto de melhoria continua.

\appendix
\section{Checklist de cumprimento do enunciado}
\begin{table}[H]
\centering
\begin{tabular}{p{0.58\textwidth}cc}
\toprule
Item de avaliacao & Projeto 1 & Projeto 2 \\
\midrule
Complemento de testes JUnit & Sim & Sim \\
Geracao automatica com EvoSuite & Sim (escopo reduzido) & Sim (escopo reduzido) \\
Cobertura estrutural & Parcial/documentada & Sim (JaCoCo) \\
Mutacao & Parcial/documentada & Sim (PIT) \\
Property-based testing & Parcial/documentada & Sim (jqwik) \\
Pipeline reproduzivel & Parcial & Sim \\
Justificacao tecnica e analise critica & Sim & Sim \\
\bottomrule
\end{tabular}
\caption{Checklist consolidada de cumprimento}
\end{table}

\section{Notas finais de interpretacao}
Os valores numericos reportados neste documento correspondem as ultimas execucoes validadas no contexto da iteracao documentada. Em iteracoes futuras, os numeros podem variar com alteracoes de codigo, de dados de teste e de configuracao de ferramentas, mantendo-se o mesmo metodo de recolha e interpretacao.

\section{Anexo Tecnico A - Fluxo detalhado de validacao}
Este anexo descreve um fluxo recomendado para repeticao controlada da avaliacao ATS.

\subsection{Passo 1 - Validacao de build minima}
\begin{itemize}
    \item Projeto 1: validar compilacao com \texttt{mvn -DskipTests compile}.
    \item Projeto 2: validar compilacao de testes com \texttt{bash ./gradlew testClasses}.
\end{itemize}

\subsection{Passo 2 - Execucao de testes base}
\begin{itemize}
    \item Projeto 1: executar testes com os perfis corretos ativos quando aplicavel.
    \item Projeto 2: executar \texttt{bash ./gradlew test} para baseline funcional.
\end{itemize}

\subsection{Passo 3 - Ferramentas de qualidade}
\begin{itemize}
    \item Projeto 2: executar cobertura e mutacao para recolher indicadores objetivos.
    \item Projeto 1: executar geracao EvoSuite no escopo reduzido definido.
\end{itemize}

\subsection{Passo 4 - Analise de resultados}
\begin{itemize}
    \item identificar classes com cobertura baixa e alto risco de regressao;
    \item identificar mutantes sobreviventes e mutantes sem cobertura;
    \item priorizar casos de teste por impacto funcional.
\end{itemize}

\subsection{Passo 5 - Iteracao de melhoria}
\begin{itemize}
    \item adicionar/ajustar testes orientados por risco e por mutantes;
    \item reexecutar pipeline completa;
    \item registar resultados no relatorio para comparacao historica.
\end{itemize}

\section{Anexo Tecnico B - Matriz de riscos operacionais}
\begin{table}[H]
\centering
\begin{tabular}{p{0.30\textwidth}p{0.14\textwidth}p{0.18\textwidth}p{0.30\textwidth}}
\toprule
Risco & Probabilidade & Impacto & Mitigacao \\
\midrule
Perfis Maven incorretos no Projeto 1 & Media & Medio/Alto & Documentar comandos por perfil e incluir comandos de referencia no relatorio \\
Runtime EvoSuite incompativel com bytecode & Alta & Medio/Alto & Fixar versoes de Java e imagem/container por ambiente controlado \\
Permissao de execucao no wrapper Gradle & Media & Medio & Usar \texttt{bash ./gradlew} ou ajustar permissao de ficheiro no repositorio \\
Mutantes sem cobertura em zonas perifericas & Alta & Medio & Criar testes orientados por sobreviventes e por no-coverage \\
Falsos positivos em sugestoes de IA & Media & Medio & Revisao manual obrigatoria e execucao local antes de aceitar alteracoes \\
\bottomrule
\end{tabular}
\caption{Matriz sintetica de riscos e mitigacoes}
\end{table}

\section{Anexo Tecnico C - Politica de qualidade para testes gerados}
Para manter consistencia e evitar inflacao artificial da suite, foi seguida a politica abaixo:
\begin{itemize}
    \item cada teste novo deve ter objetivo claro e mensuravel;
    \item testes duplicados ou equivalentes sao removidos;
    \item testes excessivamente acoplados a implementacao interna sao evitados;
    \item preferencia por nomes de teste descritivos e focados em comportamento observavel;
    \item falhas intermitentes (flaky tests) devem ser corrigidas ou removidas;
    \item toda alteracao de teste deve ser validada por execucao limpa da pipeline relevante.
\end{itemize}

\section{Anexo Tecnico D - Evidencias recomendadas para submissao}
Para uma submissao robusta e auditavel, recomenda-se anexar:
\begin{itemize}
    \item resumo de execucao JUnit dos dois projetos;
    \item relatorio HTML de JaCoCo do Projeto 2;
    \item relatorio HTML/XML de PIT do Projeto 2;
    \item logs de execucao EvoSuite para as classes alvo definidas;
    \item comandos executados e respetivos codigos de saida (exit code);
    \item versoes de Java, Maven, Gradle e Docker usadas no ambiente.
\end{itemize}

Esta lista reduz ambiguidade na avaliacao e facilita reproducao independente por parte do docente.

\end{document}
